\uuid{w8BX}
\exo7id{536}
\titre{exo7 536}
\auteur{monthub}
\organisation{exo7}
\datecreate{2001-11-01}
\isIndication{false}
\isCorrection{false}
\chapitre{Suite}
\sousChapitre{Suite définie par une relation de récurrence}
\cpge{mpsi-mp2i.nombres-reels-suites.c, pcsi.nombres-reels-suites.b, ptsi.nombres-reels-suites.b, mpsi-mp2i.nombres-reels-suites.d, pcsi.nombres-reels-suites.c, ptsi.nombres-reels-suites.c, mpsi-mp2i.nombres-reels-suites.e, pcsi.nombres-reels-suites.d, ptsi.nombres-reels-suites.d, bcpst1.suites-reelles}
\module{Analyse}
\niveau{L1}
\difficulte{}

\contenu{
\texte{
Soit $(u_n)$ la suite réelle définie par récurrence en posant $u_0=1$ et $u_{n+1}=\sqrt{1+u_n}$ si $n\in\N^{*}$.
}
\begin{enumerate}
    \item \question{Montrer que $(u_n)$ est croissante et majorée.}
    \item \question{Montrer que $(u_n)$ converge vers le nombre réel positif $\ell$ qui

vérifie $\ell^2-\ell-1=0$ et calculer $\ell$.}
\end{enumerate}
}
